\newpage

\appendix

\section{Prompt Templates}
\label{sec:prompts}

\setcounter{figure}{0}
\renewcommand{\thefigure}{A\arabic{figure}}

\begin{figure}[hbt]
\begin{mybox}[Phase 1 Prompt Template]
\begin{lstlisting}
Read the following movie review and rate its sentiment on a scale from 0 to 1, where:
  * 0 = completely negative
  * 1 = completely positive

Review:
(*@\textcolor{blue}{[REVIEW\_TEXT]}@*)

Provide ONLY your rating as a single decimal number between 0 and 1.

Examples of valid responses:
0.73
0.5
0.82

Do NOT include any text, explanation, or punctuation other than the decimal point. Your response must be ONLY the number.
\end{lstlisting}
\end{mybox}
\caption{Phase 1 prompt template for independent rating. Each model receives this prompt with an identical review text.}
\label{fig:phase1_prompt}
\end{figure}
\begin{figure}[htbp]
\begin{mybox}[Phase 2 Prompt Template (with status information)]
\begin{lstlisting}
You are now being paired with a partner for this sentiment rating task.

Your partner is (*@\textcolor{blue}{[GENDER]}@*) and has the following credentials and background:
  * Education: (*@\textcolor{blue}{[EDUCATION]}@*)
  * Experience: (*@\textcolor{blue}{[EXPERIENCE]}@*)
  * Achievement: (*@\textcolor{blue}{[ACHIEVEMENT]}@*)
  * Institution: (*@\textcolor{blue}{[INSTITUTION]}@*)
  * Expertise: (*@\textcolor{blue}{[EXPERTISE]}@*)

You will be working together on evaluating movie reviews. Your partner will also see your credentials and background.

Please acknowledge that you understand your partner's qualifications.
\end{lstlisting}
\end{mybox}

\vspace{0.5cm}

\begin{mybox}[Phase 2 Prompt Template (no status condition)]
\begin{lstlisting}
You are now being paired with a partner for this sentiment rating task.

You will be working together on evaluating movie reviews.

Please acknowledge that you understand you are working with a partner.
\end{lstlisting}
\end{mybox}
\caption{Phase 2 prompt templates for partner introduction. The top template is used in conditions with status manipulation; the bottom template is used in no-status control conditions.}
\label{fig:phase2_prompt}
\end{figure}
\begin{figure}[htbp]
\begin{mybox}[Phase 3 Prompt Template]
\begin{lstlisting}
Your rating: (*@\textcolor{blue}{[OWN\_RATING]}@*)
Your partner's rating: (*@\textcolor{blue}{[PARTNER\_RATING]}@*)

You will now have the opportunity to maintain or revise your rating.
\end{lstlisting}
\end{mybox}
\caption{Phase 3 prompt template for rating revelation. Each model sees this information neutrally presented.}
\label{fig:phase3_prompt}
\end{figure}
%\begin{figure}[htbp]
\begin{mybox}[Phase 4 Prompt Template]
\begin{lstlisting}
You will now make your final rating. Taking your partner's perspective and judgment into account is both legitimate and valuable. The most important goal is to arrive at the most accurate rating possible.

Whether your initial and final ratings match is unimportant - what matters is the accuracy of your final rating. Many raters find it helpful to consider their partner's judgment when making their final rating.

Provide your final rating as a single decimal number between 0 and 1.

Examples of valid responses:

0.73
0.5
0.82

Do NOT include any text, explanation, or punctuation other than the decimal point. Your response must be ONLY the number.
\end{lstlisting}
\end{mybox}
\caption{Phase 4 prompt template for final rating adjustment. The collective orientation framing follows \citet{Berger-1972-Status} in legitimizing consideration of partner judgment.}
\label{fig:phase4_prompt}
\end{figure}

\newpage
\section{Berger's Original Design}
\label{sec:berger}

\citet[pp. 248-250]{Berger-1972-Status} describes the experimental task as follows:

``The experiment was broken into two phases, a manipulation-phase and a decision-making phase. In the manipulation phase, subjects are placed in an expectation-state, either by testing their ability directly or by giving them status information. The actual status of the subjects is always the same. The subjects are kept apart during the experiment, for complete control over their status information. If both have the state $D_x$, each is told that one is $D_x$ but the other has some other state of $D$. For example, if both are Air Force enlisted, each is told that one is enlisted and the other is an officer; or that one is an officer and the other is enlisted. Each subject assumes that his counterpart has the other state of $D$.

In the decision phase, pairs of subjects repeat $n$ identical decision-making trials, each of which requires a binary choice. The choice is made in three stages, the first of which is the subject's initial choice between alternatives made independently, without knowing the partner's choice. Subjects then communicate these choices to their partner, after which each makes a final choice.

Subjects are instructed to make what they feel is the correct preliminary choice and, after taking the other subject's choice into consideration, to make what they feel is the correct final choice. They are told repeatedly that whether initial choices coincide with final choices is unimportant, that using others' advice is both legitimate and necessary, and that the correct final choice is of prime importance. Important to note that this method is careful to operationalized collective orientation. For example, sometimes the instructions emphasize the legitimacy of taking advice, sometimes subjects are scored as a team, and sometimes the two have been combined. 

The task the subjects complete consists of a sequence of almost identical large rectangles made up of smaller black and white rectangles. The subject must decide whether each rectangle contains more white or black area. The task is ambiguous, for the probability of a white-response for each stimulus is close to 50\%, and the decision on any given trial is independent of decisions on preceding trials. However, the subject is told that there is always a correct response, and success is defined by a set of standards giving scores (number of correct responses) typical of subjects like himself. The experimenter describes the task ability as ``contrast sensitivity'' or ``spacial'' judgment ability, and tells the subject that it is not related to artistic or mathematical ability. In other words, the experimenter attempts to preclude the subject's having prior expectations about the ability.

A precise measure of the subject's power and prestige position is obtained by studying the probability that one subject influences or is influenced by another. if the subject changes their final choice, they are said to make an $0-response$; if not, they are said to make an $S-$ or $stay$, response. The probability of an S-response measures the exercise of influence in the situation. If the theory holds, then the probability of an S-response should be greater for subjects who believe their partner is enlisted than for those who believe their partner is an officer.''